% Unit 103 terjemahan Bahasa Indonesia dari chapter7.tex baris 2679--2827.
% Sumber: Methods of Algebra, Volume 2, commit
% 9a5803ff77dd3257484cb177f851a73770a59dd3 (CC BY 4.0).
% stable-unit-id: o014.aljabr2.chapter7.exercises-a
% Cakupan: awal Latihan Bab 7 sampai latihan funktor-dg Hom;
% dilanjutkan oleh Unit 104 pada baris 2828.

% segment-id: o014.li2.chapter7.exercises-a.ex000
\begin{Exercises}
	% segment-id: o014.li2.chapter7.exercises-a.ex001
	\item Lengkapi argumen yang dihilangkan dalam
	Definisi~\ref{def:opposite-braided-algebra}, lalu buktikan bahwa untuk
	aljabar $A$ di atas kategori monoidal simetris $\mathcal{V}$, modul-$A$
	kiri sama artinya dengan modul-$A^{\opp}$ kanan.

	% segment-id: o014.li2.chapter7.exercises-a.ex002
	\item Misalkan $\mathcal{V}$ kategori monoidal beranyam dan
	$(A,\mu_A,\eta_A)$ aljabar di atasnya. Buktikan bahwa jika terdapat
	morfisme $\mu':A\otimes A\to A$ dan $\eta':\munit\to A$ dalam
	$\cate{Alg}(\mathcal{V})$ yang membuat $A$ menjadi aljabar di atas
	$\cate{Alg}(\mathcal{V})$, maka niscaya $\eta'=\eta_A$ dan
	$\mu'=\mu_A$. Khususnya, dengan menerapkan
	Lema~\ref{prop:CAlg-mu}, diperoleh bahwa $A$ komutatif. Dengan demikian,
	buktikan \eqref{eqn:AlgAlg-CAlg}.

	% segment-id: o014.li2.chapter7.exercises-a.hi001
	\begin{hint}
		Mudah ditunjukkan bahwa $\eta'=\eta_A$. Untuk $\mu'=\mu_A$, tuliskan
		$\mu_{A\otimes A}$ bagi perkalian pada $A\otimes A$. Syarat yang
		diberikan menghasilkan diagram komutatif
		% segment-id: o014.li2.chapter7.exercises-a.q001
		\[\begin{tikzcd}[column sep=large]
			(A \otimes A) \otimes (A \otimes A) \arrow[r, "{\mu' \otimes \mu'}"] \arrow[d, "{\mu_{A \otimes A}}"'] & A \otimes A \arrow[d, "{\mu_A}"] \\
			A \otimes A \arrow[r, "{\mu'}"'] & A.
		\end{tikzcd}\]
		Jika empat salinan $A$ ditempatkan dalam tabel $2\times2$, diagram
		komutatif tersebut dapat dibayangkan sebagai hukum pertukaran horizontal
		dan vertikal
		% segment-id: o014.li2.chapter7.exercises-a.g001
		\begin{center}\begin{tikzpicture}[baseline=(O)]
			\matrix (M) [matrix of nodes] {
				$A$ & $A$ \\
				$A$ & $A$ \\
			};
			\coordinate (O) at (M-1-1.south west);
			\draw (M-1-1.north west) rectangle (M-1-2.south east);
			\draw (M-2-1.north west) rectangle (M-2-2.south east);
		\end{tikzpicture} = \begin{tikzpicture}[baseline=(O)]
			\matrix (M) [matrix of nodes] {
				$A$ & $A$ \\
				$A$ & $A$ \\
			};
			\coordinate (O) at (M-1-1.south west);
			\draw (M-1-1.north west) rectangle (M-2-1.south east);
			\draw (M-1-2.north west) rectangle (M-2-2.south east);
		\end{tikzpicture} \quad \begin{tabular}{rl}
			Horizontal: & kalikan dengan $\mu'$ \\
			Vertikal: & kalikan dengan $\mu_A$
		\end{tabular}\end{center}
		Sekarang argumen untuk $\mu'=\mu_A$ dapat digambarkan sebagai
		% segment-id: o014.li2.chapter7.exercises-a.g002
		\begin{center}\begin{tikzpicture}[baseline=(O)]
			\matrix (M) [matrix of nodes] {
				$A$ & $A$ \\
			};
			\coordinate (O) at (M-1-1.south west);
			\draw (M-1-1.north west) rectangle (M-1-2.south east);
		\end{tikzpicture}
		$=$ \begin{tikzpicture}[baseline=(O)]
			\matrix (M) [matrix of nodes, every node/.style={anchor=base, minimum width=1.7em, minimum height=1.7em}] {
				$A$ & $\munit$ \\
				$\munit$ & $A$ \\
			};
			\coordinate (O) at (M-1-1.south west);
			\draw (M-1-1.north west) rectangle (M-2-1.south east);
			\draw (M-1-2.north west) rectangle (M-2-2.south east);
		\end{tikzpicture}
		$=$ \begin{tikzpicture}[baseline=(O)]
			\matrix (M) [matrix of nodes, every node/.style={anchor=base, minimum width=1.5em, minimum height=1.7em}] {
				$A$ & $\munit$ \\
				$\munit$ & $A$ \\
			};
			\coordinate (O) at (M-1-1.south west);
			\draw (M-1-1.north west) rectangle (M-1-2.south east);
			\draw (M-2-1.north west) rectangle (M-2-2.south east);
		\end{tikzpicture}
		$=$ \begin{tikzpicture}[baseline=(O)]
			\matrix (M) [matrix of nodes] {
				$A$ \\
				$A$ \\
			};
			\coordinate (O) at (M-1-1.south west);
			\draw (M-1-1.north west) rectangle (M-2-1.south east);
		\end{tikzpicture}\end{center}
		Cobalah juga menerjemahkannya menjadi diagram komutatif atau persamaan
		yang bersesuaian.
	\end{hint}

	% segment-id: o014.li2.chapter7.exercises-a.ex003
	\item Tinjau aljabar-dg $A$ di atas $\mathcal{A}$
	(Definisi~\ref{def:dg-algebra}). Periksa bahwa bila $A=\munit$ (objek
	satuan), konsep modul-$A$ kiri maupun modul-$A$ kanan sama dengan konsep
	kompleks di atas $\mathcal{A}$.

	% segment-id: o014.li2.chapter7.exercises-a.ex004
	\item Tetapkan gelanggang komutatif $\Bbbk$. Misalkan $A$ aljabar-dg di
	atas $\Bbbk\dcate{Mod}$ dan $M$ kompleks. Buktikan bahwa memberikan
	struktur modul-$A$ kiri pada $M$ sama artinya dengan memberikan
	homomorfisme aljabar-dg
	$A\to\End^\bullet(M):=\Hom^\bullet(M,M)$.

	% segment-id: o014.li2.chapter7.exercises-a.ex005
	\item Misalkan $\Bbbk$ gelanggang komutatif dan $R$ aljabar-dg di atas
	$\Bbbk\dcate{Mod}$. Untuk modul-$R$ kanan $X$ dan modul-$R$ kiri $Y$,
	kita masing-masing mengidentifikasi $R$ dan $X,Y$ dengan aljabar-$\Bbbk$
	dan modul-$\Bbbk$ yang dilengkapi dekomposisi jumlah langsung
	% segment-id: o014.li2.chapter7.exercises-a.q002
	\[ R = \bigoplus_n R^n, \quad X = \bigoplus_n X^n, \quad Y = \bigoplus_n Y^n \]
	serta endomorfisme berderajat $1$, $d$, yang memenuhi $d^2=0$. Karena itu,
	$X\dotimes{R}Y$ sudah terdefinisi sebagai modul-$\Bbbk$.
	\begin{enumerate}[(i)]
		\item Ingat bahwa
		$X\otimes Y=\bigoplus_n(\bigoplus_{p+q=n}X^p\otimes Y^q)$ sekaligus
		merupakan modul-$\Bbbk$ dan mempunyai struktur kompleks. Untuk
		homomorfisme modul-$\Bbbk$ yang jelas
		$\pi:X\otimes Y\twoheadrightarrow X\dotimes{R}Y$, periksa bahwa
		$\Ker(\pi)$ adalah subkompleks, sehingga $X\dotimes{R}Y$ mempunyai
		struktur kompleks kanonik.

		\item Buktikan bahwa jika $X$ bimodul-$(A,R)$ dan $Y$ bimodul-$(R,B)$,
		dengan $A$ dan $B$ juga aljabar-dg, maka $X\dotimes{R}Y$ secara alami
		menjadi bimodul-$(A,B)$. Gunakan hal ini untuk memperluas adjoin dalam
		Proposisi~\ref{prop:tensor-Hom-cplx} ke kasus modul-dg. Kompleks
		$\Hom^\bullet(Y_B,Z_B)$ dan $\Hom^\bullet({}_AX,{}_AZ)$ yang muncul di
		sini didefinisikan seperti dalam Contoh~\ref{eg:dg-Hom-cplx}, tetapi
		tanda-tanda perlu diperhatikan ketika memberikan struktur bimodul-
		$(A,R)$ dan bimodul-$(R,B)$.
	\end{enumerate}

	% segment-id: o014.li2.chapter7.exercises-a.ex006
	\item Tetapkan aljabar-dg $A$ di atas $\mathcal{A}$. Jika
	$p>0\implies A^p=0$, maka $A$ disebut \emph{nonpositif}. Untuk modul-$A$
	kiri (atau kanan) $M$ dan $n\in\Z$, ambil subkompleks terpotong
	$\tau^{\leq n}M$ menurut Definisi~\ref{def:truncation-cplx}. Buktikan
	bahwa bila $A$ nonpositif:
	\begin{enumerate}[(i)]
		\item setiap $d_M^p:M^p\to M^{p+1}$ bersifat $A^0$-linear;
		\item terdapat barisan eksak pendek modul-$A$
		$0\to\tau^{\leq n}M\to M\to\tilde\tau^{\geq n+1}M\to0$.
	\end{enumerate}
	\index{aljabar dg!tak-positif@tak-positif (non-positive)}

	% segment-id: o014.li2.chapter7.exercises-a.ex007
	\item Tetapkan aljabar-dg $A$ di atas $\mathcal{A}$. Ingat bahwa semua
	modul-$A$ kiri membentuk kategori-dg $A\dcate{dgMod}$
	(Contoh~\ref{eg:dg-Hom-cplx}); Catatan~\ref{rem:dg-homotopy-cat} kemudian
	memberikan kategori homotopi $\mathrm{h}(A\dcate{dgMod})$. Modul-$A$ kiri
	$M$ disebut asiklik bila ia asiklik sebagai kompleks; morfisme modul-$A$
	$f:M\to N$ disebut kuasi-isomorfisme bila ia kuasi-isomorfisme sebagai
	morfisme kompleks.
	\begin{enumerate}[(i)]
		\item Misalkan $M$ modul-$A$ kiri yang strukturnya ditentukan oleh
		keluarga morfisme dalam $\mathcal{A}$,
		$\alpha_{A,M}^{p,q}:A^p\otimes M^q\to M^{p+q}$ ($p,q\in\Z$).
		Tunjukkan bahwa rumus berikut memberikan struktur modul-$A$ kiri pada
		kompleks $M[1]$:
		% segment-id: o014.li2.chapter7.exercises-a.q003
		\[ \alpha_{A, M[1]}^{p, q} := (-1)^p \alpha_{A, M}^{p, q+1} . \]
		Selanjutnya, tunjukkan bahwa $M\mapsto M[1]$ memberikan funktor-dg dari
		$A\dcate{dgMod}$ ke dirinya sendiri.
		\item Untuk sembarang morfisme $f:M\to N$ dalam $A\dcate{dgMod}$,
		tunjukkan bahwa kerucut pemetaan $\Cone(f)$ memiliki struktur modul-$A$
		kiri kanonik, yang dalam notasi di atas diberikan oleh
		% segment-id: o014.li2.chapter7.exercises-a.q004
		\[ \alpha_{A, \Cone(f)}^{p, q} = \alpha_{A, M[1]}^{p, q} \oplus \alpha_{A, N}^{p, q}. \]
		\item Mengikuti konstruksi dalam \S\ref{sec:derived-cat}, gunakan kerucut
		pemetaan untuk membuat $\mathrm{h}(A\dcate{dgMod})$ menjadi kategori
		triangulasi. Selanjutnya, lakukan pelokalan Verdier pada modul-$A$
		asiklik, atau, dengan kata lain, tambahkan invers kuasi-isomorfisme,
		untuk memperoleh kategori turunan $\cate{D}(A\dcate{dgMod})$.
		Definisikan $\cate{D}^{\star}(A\dcate{dgMod})$ yang bersesuaian untuk
		$\star\in\{+,-,\bdd\}$.
		\item Dengan asumsi $A$ nonpositif, perluas
		Proposisi~\ref{prop:derived-full-faithful-subcat} ke
		$\cate{D}(A\dcate{dgMod})$ dan
		$\cate{D}^{\star}(A\dcate{dgMod})$, untuk
		$\star\in\{+,-,\bdd\}$.
		% segment-id: o014.li2.chapter7.exercises-a.hi002
		\begin{hint}
			Gunakan funktor pemotongan.
		\end{hint}
		\item Dengan mengacu pada \S\ref{sec:K-injectives}, definisikan
		modul-$A$ K-injektif dan K-projektif, lalu gunakan objek-objek ini untuk
		mendeskripsikan morfisme dalam $\cate{D}(A\dcate{dgMod})$. Tempatkan
		semuanya dalam kerangka \S\ref{sec:triangulated-functor-localization}
		untuk membahas hubungannya dengan funktor turunan.
		\item Untuk modul-$A$ kanan $M$, tunjukkan bahwa kompleks $M[1]$
		mempunyai struktur modul-$A$ kanan
		$\alpha_{M[1],A}^{p,q}=\alpha_{M,A}^{p,q+1}$ dan bahwa semua sifat di
		atas tetap berlaku.
	\end{enumerate}

	% segment-id: o014.li2.chapter7.exercises-a.ex008
	\item Misalkan $\mathcal{C}$ kategori-dg di atas $\Bbbk$. Definisikan
	$\mathcal{C}^{\opp}$ melalui
	$\Obj(\mathcal{C}^{\opp})=\Obj(\mathcal{C})$ dan, untuk sembarang
	$X,Y\in\Obj(\mathcal{C}^{\opp})$,
	$\Hom^\bullet_{\mathcal{C}^{\opp}}(X,Y):=
	\Hom^\bullet_{\mathcal{C}}(Y,X)$. Definisikan komposisinya sebagai
	% segment-id: o014.li2.chapter7.exercises-a.q005
	\begin{align*}
		\Hom^p_{\mathcal{C}^{\opp}}(Y, Z) \dotimes{\Bbbk} \Hom^q_{\mathcal{C}^{\opp}}(X, Y) & \to \Hom^{p+q}_{\mathcal{C}^{\opp}}(X, Z) \quad p, q \in \Z \\
		f \otimes g & \mapsto (-1)^{pq} \underbracket{f \circ g}_{\text{operasi dalam }\Hom^\bullet_{\mathcal{C}}}.
	\end{align*}
	Periksa bahwa ini memang memberikan kategori-dg $\mathcal{C}^{\opp}$,
	yang disebut kategori-dg lawan dari $\mathcal{C}$.

	% segment-id: o014.li2.chapter7.exercises-a.ex009
	\item Tetapkan gelanggang komutatif $\Bbbk$, aljabar-dg $R$ di atas
	$\Bbbk\dcate{Mod}$, modul-$R$ kanan $X$, dan modul-$R$ kiri $Y$. Buktikan
	bahwa funktor yang diperkenalkan sebelumnya,
	% segment-id: o014.li2.chapter7.exercises-a.q006
	\[ X \dotimes{R} (\cdot): R\dcate{dgMod} \to \cate{C}(\Bbbk), \quad (\cdot) \dotimes{R} Y: \cated{dgMod}R \to \cate{C}(\Bbbk), \]
	secara alami diangkat menjadi funktor-dg. Rumuskan garis besar hasil yang bersesuaian
	ketika $X$ atau $Y$ juga mempunyai struktur bimodul.

	% segment-id: o014.li2.chapter7.exercises-a.hi003
	\begin{hint}
		Kuncinya ialah mendefinisikan morfisme antarkompleks $\Hom$. Untuk
		$X\dotimes{R}(\mathord\cdot)$, misalkan $M$ dan $N$ modul-$R$ kiri.
		Aturan tanda Koszul mensyaratkan citra
		$f\in\Hom^n_R(M,N)$, yaitu
		$\identity_X\otimes f\in
		\Hom^n(X\dotimes{R}M,X\dotimes{R}N)$, didefinisikan oleh
		% segment-id: o014.li2.chapter7.exercises-a.q007
		\[ (\identity_X \otimes f)(x \otimes m) = (-1)^{pn} x \otimes f(m), \quad x \in X^p, \; m \in M^q. \]
		Perlu diperiksa bahwa
		$\identity_X\otimes d_{\Hom^\bullet}(f)=
		d_{\Hom^\bullet}(\identity_X\otimes f)$.

		Kasus $(\mathord\cdot)\dotimes{R}Y$ lebih sederhana: cukup ambil
		$(f\otimes\identity_Y)(m\otimes y)=f(m)\otimes y$.
	\end{hint}

	% segment-id: o014.li2.chapter7.exercises-a.ex010
	\item Misalkan $\mathcal{C}$ kategori-dg di atas gelanggang komutatif
	$\Bbbk$ dan $M\in\Obj(\mathcal{C})$.
	\begin{enumerate}[(i)]
		\item Angkat
		$\Hom^\bullet(M,\mathord\cdot):\mathcal{C}\to\cate{C}(\Bbbk)$ menjadi
		funktor-dg.

		% segment-id: o014.li2.chapter7.exercises-a.hi004
		\begin{hint}
			Misalkan $X,Y\in\Obj(\mathcal{C})$ dan
			$f\in\Hom^n(X,Y)$. Tetapkan
			$\mathcal{H}_X=\Hom^\bullet(M,X)$. Citra $f$ harus diambil sebagai
			% segment-id: o014.li2.chapter7.exercises-a.q008
			\[ \left[ f_*: g \xmapsto{\text{morfisme berderajat }n} fg \right] \in \Hom^n(\mathcal{H}_X, \mathcal{H}_Y). \]
			Gunakan definisi kompleks $\Hom$ untuk memeriksa bahwa
			$(d_{\Hom^\bullet(X,Y)}f)_*=
			d_{\Hom^\bullet(\mathcal{H}_X,\mathcal{H}_Y)}(f_*)$.
		\end{hint}
		\item Angkat
		$\Hom^\bullet(\mathord\cdot,M):\mathcal{C}\to
		\cate{C}(\Bbbk)^{\opp}$ menjadi funktor-dg.

		% segment-id: o014.li2.chapter7.exercises-a.hi005
		\begin{hint}
			Tetapkan $\mathcal{K}_X=\Hom^\bullet(X,M)$. Aturan tanda Koszul
			mensyaratkan citra $f\in\Hom^n(X,Y)$ diambil sebagai
			% segment-id: o014.li2.chapter7.exercises-a.q009
			\[ \left[\begin{array}{c}
				f^*: g \xmapsto{\text{morfisme berderajat }n} (-1)^{np} gf \\
				\forall p \in \Z , \; \forall\; g \in \Hom^p(Y, M)
			\end{array}\right] \in \Hom^n(\mathcal{K}_Y, \mathcal{K}_X). \]
			Periksa dengan cermat bahwa
			\begin{compactitem}
				\item $(d_{\Hom^\bullet(X,Y)}f)^*=
				d_{\Hom^\bullet(\mathcal{K}_Y,\mathcal{K}_X)}(f^*)$;
				\item $(fh)^*$ sama dengan komposisi $h^*$ dan $f^*$ dalam
				$\Hom^\bullet_{\cate{C}(\Bbbk)}$.
			\end{compactitem}
		\end{hint}
	\end{enumerate}
